\chapter{Asking LeanFM}\label{ch:mcp}

Everything in Parts~\ref{part:language} and~\ref{part:ctl} is available to an agent
through LeanFM's server, which speaks the Model Context Protocol (MCP). This chapter
connects Claude to it and then works through the questions of the book as an agent asks
them: each section shows the request the agent sends and an excerpt of the answer
LeanFM returns, exactly as generated. The answers are shortened only where a string or
list is long; nothing is rewritten.

\section{Connecting}

The server runs locally. For Claude Desktop it runs over standard input and output, so
it needs no port and no key:
\begin{lstlisting}[language={},basicstyle=\ttfamily\footnotesize]
lake build leanfm-mcp
python3 scripts/claude_desktop_config.py --root "$PWD" --install
\end{lstlisting}
The script adds a \texttt{leanfm} entry to Claude Desktop's configuration that starts
\texttt{leanfm-mcp} in the repository. After a restart, Claude lists the tools below and
reads the language references the server publishes as resources:
\begin{center}
\texttt{leanfm://language/session-reference}\\
\texttt{leanfm://language/logic-reference}
\end{center}
The same tools are served over HTTP by \texttt{make serve}, at
\texttt{http://127.0.0.1:8080/mcp}, for clients that connect by URL.

A conversation with an agent then has a fixed shape. The person describes what the
software must do; the agent writes a protocol and checks it; LeanFM answers with
verdicts, paths and diagrams; the person pushes back; the protocol changes. Only after
the protocol says what was meant is code written, and the code is then checked against
the protocol's logs.

\section{The tools}

{\footnotesize\begin{tabular}{@{}p{.25\textwidth}p{.33\textwidth}p{.36\textwidth}@{}}
\toprule
Tool & What it does & Arguments (* required) \\
\midrule
\texttt{validate\_\allowbreak{}generated\_\allowbreak{}artifacts} & Validate generated artifacts & \texttt{limits} \\
\texttt{requirements\_\allowbreak{}interrogation} & Read requirements interrogation & \texttt{limits} \\
\texttt{code\_\allowbreak{}generation\_\allowbreak{}prompt} & Read one-shot implementation prompt & \texttt{limits} \\
\texttt{explain\_\allowbreak{}formula} & Explain a CTL formula & \texttt{formula*, limits, sentence, task} \\
\texttt{check\_\allowbreak{}ctl} & Check a CTL formula against a task & \texttt{formula*, limits, task*} \\
\texttt{check\_\allowbreak{}trace} & Check an event trace and draw it & \texttt{drafts, events*, limits} \\
\texttt{classify\_\allowbreak{}requirement} & Classify requirement observability & \texttt{formula, limits, predicates, task} \\
\texttt{check\_\allowbreak{}protocol} & Find attacks on a security protocol & \texttt{attacker*, defender, domain, domains, formula, limits, roles*} \\
\texttt{check\_\allowbreak{}reassembly} & Find stream-reassembly evasions & \texttt{box, formula, host, limits, mode, signature} \\
\texttt{knows} & What an actor can derive & \texttt{actor*, at*, events*, initial, limits, query, seesAll} \\
\texttt{session\_\allowbreak{}check} & Check a session type & \texttt{example, file, lean, limits, mayCall, notHttp, overlay, project, proto, protocol} \\
\texttt{session\_\allowbreak{}diagram} & Draw a session type & \texttt{example, file, lean, limits, overlay, project, proto, protocol, theme} \\
\texttt{session\_\allowbreak{}mdp} & Render a session decision process and terminal probabilities & \texttt{events, example, file, horizon, lean, limits, optimize, overlay, project, proto, protocol, theme, view, walk} \\
\texttt{session\_\allowbreak{}sla} & Check a session type's service-level requirements & \texttt{at, events, example, file, lean*, limits, overlay, project, proto, theme} \\
\texttt{session\_\allowbreak{}simulate} & Generate an event stream from a session type & \texttt{example, file, instances, lean, limits, loops, maxEvents, overlay, project, proto, protocol, seed, spread} \\
\texttt{session\_\allowbreak{}conform} & Check a recorded run against a session type & \texttt{allowIncomplete, events*, example, file, horizonMs, lean, limits, overlay, project, proto, protocol} \\
\texttt{session\_\allowbreak{}spans} & Time the calls in a log, without trusting clocks across machines & \texttt{events*, example, file, lean, limits, project, proto, protocol} \\
\texttt{session\_\allowbreak{}next} & What may happen next in a session & \texttt{events, example, file, lean, limits, overlay, project, proto, protocol} \\
\texttt{session\_\allowbreak{}ctl} & Check a temporal property of a session type & \texttt{example, faithful, file, formula, lean, limits, overlay, project, proto, protocol} \\
\texttt{session\_\allowbreak{}instances} & Find deadlocks and races between instances of a system & \texttt{example, file, instances, lean*, limits, perInstance, project, properties, proto} \\
\texttt{session\_\allowbreak{}realizable} & Check what a deployment can actually do & \texttt{example, file, lean*, limits, project} \\
\texttt{session\_\allowbreak{}deployment} & Check where a protocol's roles run & \texttt{deployment, example, file, lean*, limits, project, theme} \\
\texttt{session\_\allowbreak{}review} & Review a session type design & \texttt{example, file, lean, limits, overlay, project, proto, protocol} \\
\texttt{session\_\allowbreak{}openapi} & A role's HTTP interface as OpenAPI & \texttt{example, file, lean, limits, overlay, project, proto, protocol, role*} \\
\texttt{session\_\allowbreak{}types} & The message types as protobuf, Lean and JSON schema & \texttt{example, file, lean, limits, project, proto, protocol} \\
\texttt{session\_\allowbreak{}compat} & Which roles can upgrade on their own & \texttt{example, file, lean, limits, old, overlay, project, proto, protocol} \\
\texttt{session\_\allowbreak{}actor} & A protocol as one role sees it & \texttt{example, file, lean, limits, overlay, project, proto, protocol, role*, theme} \\
\texttt{session\_\allowbreak{}dataset} & Generate many sessions and summarize them & \texttt{example, file, lean, limits, overlay, project, proto, protocol, seed, sessions} \\
\texttt{session\_\allowbreak{}explore} & Browse a protocol as a tree & \texttt{example, file, focus, lean, limits, project, proto, states, theme} \\
\texttt{session\_\allowbreak{}skeleton} & Generate code to implement a session type & \texttt{example, file, language, lean, limits, project, proto, protocol} \\
\bottomrule
\end{tabular}}


The session tools take a protocol as \texttt{lean} (LeanFM protocol text), as
\texttt{example} (a protocol the server carries, such as \texttt{goserver}) or as
\texttt{file} (a protocol file of the project). The logic tools take a \texttt{task}:
a canonical id or an inline draft of states and transitions.

\section{Is the protocol playable?}

The first call on any new protocol is \texttt{session\_check}: projection and
well-formedness (Chapter~\ref{ch:projection}).

\begin{lstlisting}[language={},basicstyle=\ttfamily\footnotesize,breaklines=true]
tools/call session_check
{"lean": "protocol Images (Browser, Server) { ... }"}

answer (excerpt)
{"verdict": "ok",
 "reading": "Images: no issues; 2 roles, 3 messages, 1 top-level calls",
 "issues": []}
\end{lstlisting}


\section{Does a property hold?}

\texttt{session\_ctl} with a \texttt{formula} decides it over the protocol's checker
model (Chapter~\ref{ch:checker}). The formula may be JSON or an s-expression. Every
image request is answered:

\begin{lstlisting}[language={},basicstyle=\ttfamily\footnotesize,breaklines=true]
tools/call session_ctl
{"lean": "protocol Images ...",
 "formula": "(AG (implies (after GetImage) (AF (or (after Image) (after NotFound)))))"}

answer (excerpt)
{"verdict": "proved",
 "states": 4,
 "report": "",
 "reading":
 "Images: on every path, always (if GetImage was just delivered then (on every path, eventually (Image was just delivered or NotFound was just delivered ...",
 "notation": "AG (GetImage → AF (Image ∨ NotFound))"}
\end{lstlisting}


A property that does not hold comes back with the path that breaks it. The server is
not obliged to answer with an image:

\begin{lstlisting}[language={},basicstyle=\ttfamily\footnotesize,breaklines=true]
tools/call session_ctl
{"lean": "protocol Images ...", "formula": "(AG (implies (after GetImage) (AX (after Image))))"}

answer (excerpt)
{"verdict": "disproved",
 "report": "counterexample: GetImage -> NotFound",
 "notation": "AG (GetImage → AX Image)"}
\end{lstlisting}


Order is a property too. No NotFound comes before the request it answers:

\begin{lstlisting}[language={},basicstyle=\ttfamily\footnotesize,breaklines=true]
tools/call session_ctl
{"lean": "protocol Images ...",
 "formula": "(not (someUntil (not (on GetImage)) (on NotFound)))"}

answer (excerpt)
{"verdict": "proved", "report": "", "notation": "¬(¬GetImage someUntil NotFound)"}
\end{lstlisting}


\section{What should I ask?}

Without a formula, \texttt{session\_ctl} builds properties from the protocol and decides
each. For PlaceOrder:

\begin{lstlisting}[language={},basicstyle=\ttfamily\footnotesize,breaklines=true]
PlaceOrder: 12 properties suggested from the protocol, 10 proved; pass one as formula to decide it with its path

proved: Every session can still end from every state it reaches (no trap)
  always possibly stop

proved: Every session ends (no path runs forever)
  inevitably stop

disproved: Some session can run forever
  someAlways ¬stop

proved: Every session starts with CreateOrder
  allNext CreateOrder

disproved: No secret is ever sent
  always ¬secret

proved: No secret ever reaches Client
  always ¬secret →Client

proved: Every CreateOrder is answered (Quote or SoldOut) before the session ends
  always (CreateOrder → (¬stop allUntil (Quote ∨ SoldOut)))

proved: CreateOrder can happen, and every CreateOrder is eventually answered
  CreateOrder ⇒ inevitably (Quote ∨ SoldOut)

proved: No Quote or SoldOut before a CreateOrder
  ¬(¬CreateOrder someUntil (Quote ∨ SoldOut))

proved: Every Pay is answered (Receipt or PaymentDeclined) before the session ends
  always (Pay → (¬stop allUntil (Receipt ∨ PaymentDeclined)))

proved: Pay can happen, and every Pay is eventually answered
  Pay ⇒ inevitably (Receipt ∨ PaymentDeclined)

proved: No Receipt or PaymentDeclined before a Pay
  ¬(¬Pay someUntil (Receipt ∨ PaymentDeclined))
\end{lstlisting}


Two are disproved, and both are informative. A purchase always ends, so ``some session
can run forever'' is false, which is what was wanted. ``No secret is ever sent'' is
false because the client sends its card number to the shop and the shop to the bank: a
project that must keep card numbers away from some role states that rule instead, for
example $\always\neg\mathsf{secret}(\textit{to}: \textit{Log})$, and LeanFM decides it the
same way. A suggestion is a question the protocol answers, not a rule it must keep.

\section{The protocol as CTL}

\texttt{faithful: "rules"} answers the typestates, the lookahead and the faithful rules
of Chapter~\ref{ch:faithful}, each as an s-expression with its shared parts labelled and,
when it is small, in CTL notation:

\begin{lstlisting}[language={},basicstyle=\ttfamily\footnotesize,breaklines=true]
tools/call session_ctl
{"lean": "protocol Images ...", "faithful": "rules"}

answer (excerpt)
{"verdict": "rules",
 "typestates": 4,
 "rules":
 [{"subformulas": 8,
   "notation": "¬(GetImage ∨ Image ∨ NotFound) ∧ ¬terminal",
   "name": "start",
   "formula":
   "(and (not (or (after GetImage) (after Image) (after NotFound))) (not terminal))"},
  {"subformulas": 15,
   "notation":
   "AG ((¬(GetImage ∨ Image ∨ NotFound) ∧ ¬terminal) → (EX (GetImage ∧ ¬terminal) ∧ AX (GetImage ∧ ¬terminal)))",
   "name": "typestate 0: (start)",
   "formula":
   "(AG (implies (and (not (or (after GetImage) (after Image) (after NotFound))) #0=(not terminal)) (and (EX #1=(and (after GetImage) #0#)) (AX (or #1#)))))"},
  {"subformulas": 15,
   "notation":
   "AG ((GetImage ∧ ¬terminal) → (EX (Image ∧ terminal) ∧ EX (NotFound ∧ terminal) ∧ AX ((Image ∧ terminal) ∨ (NotFound ∧ terminal))))",
   "name": "typestate 1: GetImage",
   "formula":
   "(AG (implies (and (after GetImage) (not terminal)) (and (EX #0=(and (after Image) terminal)) (EX #1=(and (after NotFound) terminal)) (AX (or #0# #1#)))))"},
  {"subformulas": 9,
   "notation": "AG ((Image ∧ terminal) → (EX (Image ∧ terminal) ∧ AX (Image ∧ terminal)))",
   "name": "typestate 2: Image (may end)",
   "formula": "(AG (implies #0=(and (after Image) terminal) (and (EX #0#) (AX (or #0#)))))"},
  {"subformulas": 9,
   "notation":
   "AG ((NotFound ∧ terminal) → (EX (NotFound ∧ terminal) ∧ AX (NotFound ∧ terminal)))",
   "name": "typestate 3: NotFound (may end)",
   "formula":
   "(AG (implies #0=(and (after NotFound) terminal) (and (EX #0#) (AX (or #0#)))))"}],
 "ruleCount": 5,
 "lookahead": 0,
 "deterministic": true}
\end{lstlisting}


\section{Is this implementation faithful?}

\texttt{faithful} with a \texttt{task} (a draft of states and transitions, for example
one reconstructed from an implementation) or \texttt{lean} (another protocol text)
decides faithfulness as one bisimulation check and, when it fails, names where:

\begin{lstlisting}[language={},basicstyle=\ttfamily\footnotesize,breaklines=true]
tools/call session_ctl
{"lean": "protocol Images ...",
 "faithful": {"lean": "protocol Images ... (no NotFound branch)"}}

answer (excerpt)
{"verdict": "unfaithful",
 "reading":
 "Images: the candidate (3 nodes) is not faithful to the protocol (4 typestates), decided by walk: after GetImage, the protocol allows Image (may end), NotFound (may end) and the candidate Image (may end)",
 "protocolNext": ["Image (may end)", "NotFound (may end)"],
 "path": ["GetImage"],
 "method": "walk",
 "candidateNext": ["Image (may end)"]}
\end{lstlisting}


\section{What may happen next?}

\texttt{session\_next} takes the events of one session so far and lists every step the
protocol allows next, with the event each must name as its cause. It is what a server
uses to decide whether to accept an incoming message.

\begin{lstlisting}[language={},basicstyle=\ttfamily\footnotesize,breaklines=true]
tools/call session_next
{"lean": "protocol PlaceOrder ...",
 "events": [{"to": "Shop", "label": "CreateOrder", "id": "e1", "from": "Client"}]}

answer (excerpt)
{"verdict": "running",
 "reading": "PlaceOrder: after 1 event, 2 steps may happen next",
 "next":
 [{"to": "Client", "label": "Quote", "from": "Shop", "cause": "e1"},
  {"to": "Client", "label": "SoldOut", "from": "Shop", "cause": "e1"}]}
\end{lstlisting}


\section{A formula with shared parts}

\texttt{check\_ctl} decides a formula over an inline draft. Written as an s-expression,
a part used twice is written once:

\begin{lstlisting}[language={},basicstyle=\ttfamily\footnotesize,breaklines=true]
tools/call check_ctl
{"task":
 {"transitions":
  [{"src": "s0", "message": "a", "dst": "s1"},
   {"src": "s1", "message": "b", "dst": "s0"},
   {"src": "s1", "message": "c", "dst": "s2"}],
  "terminals": ["s2"],
  "states": ["s0", "s1", "s2"],
  "initial": "s0",
  "id": "loopy"},
 "formula": "(and #1=(EF (after c)) (AG (implies (after a) #1#)))"}

answer (excerpt)
{"verdict": "proved",
 "reading":
 "(it is possible that eventually c was just delivered) and (on every path, always (if a was just delivered then (it is possible that eventually c was j ...",
 "reachableNodes": 4,
 "notation": "EF c ∧ AG (a → EF c)"}
\end{lstlisting}


\section{Does this log follow the protocol?}

\texttt{session\_conform} replays a log, derives its sessions when the log does not name
them, and checks the where clauses on the values (Chapter~\ref{ch:where}):

\begin{lstlisting}[language={},basicstyle=\ttfamily\footnotesize,breaklines=true]
tools/call session_conform
{"lean": "protocol Images ...",
 "events":
 [{"to": "Server", "label": "GetImage", "id": "a1", "from": "Browser", "fields": {"id": "cat"}},
  {"to": "Server", "label": "GetImage", "id": "b1", "from": "Browser", "fields": {"id": "dog"}},
  {"to": "Browser", "label": "NotFound", "id": "b2", "from": "Server", "cause": "b1"},
  {"to": "Browser",
   "label": "Image",
   "id": "a2",
   "from": "Server",
   "fields": {"id": "cat", "bytes": "..."},
   "cause": "a1"}]}

answer (excerpt)
{"verdict": "conforms",
 "reading": "Images: 4 events in 2 sessions follow the protocol",
 "derivedSessions":
 {"unattributed": [],
  "sessions":
  [{"key": "chain from a1", "events": [0, 3], "count": 2},
   {"key": "chain from b1", "events": [1, 2], "count": 2}],
  "reading":
  "2 sessions derived from causes and the declared session key; the log names no sessions",
  "omitted": 0}}
\end{lstlisting}


\section{What interface does it imply?}

\texttt{session\_openapi} writes the OpenAPI document of one role
(Section~\ref{sec:openapi} shows the whole document for Images); requests without a
route are listed as unrouted rather than guessed:

\begin{lstlisting}[language={},basicstyle=\ttfamily\footnotesize,breaklines=true]
paths:
  "/images/{id}":
    get:
      operationId: "GetImage"
      parameters:
        - in: "query"
          name: "id"
          required: true
          schema:
            type: "string"
      responses:
        "200":
          content:
            "application/json":
              schema:
                "$ref": "#/components/schemas/Image"
          description: "Image"
        "404":
          content:
            "application/json":
              schema:
                "$ref": "#/components/schemas/NotFound"
          description: "NotFound"
      summary: "GetImage from Browser"
components:
  schemas:
    GetImage:
      additionalProperties: false
      properties:
        id:
\end{lstlisting}


\section{When an answer stops at a limit}

Every bound that can stop LeanFM's work is a named limit: \texttt{states} explored,
\texttt{logEvents} in a log, \texttt{formulaNodes} in a formula, and the others the
server's instructions list with their defaults. A limit is a warning, not a wall. An
answer that stops at one names it and says how to raise it; here the limit is set
absurdly low to show the refusal:

\begin{lstlisting}[language={},basicstyle=\ttfamily\footnotesize,breaklines=true]
tools/call session_ctl
{"limits": {"states": 5}, "lean": "protocol PlaceOrder ...", "formula": {"inevitably": "stop"}}

answer (excerpt)
{"verdict": "rejected",
 "reason":
 "more than 5 protocol states; the property was not checked; this is the states limit (5); to allow more, call again with \"limits\": {\"states\": 10} (a warning, not a wall: raise it when the machine has the memory)"}
\end{lstlisting}


Every tool takes the optional \texttt{limits} argument, for that call only. With the
limit raised, the answer says what was used:

\begin{lstlisting}[language={},basicstyle=\ttfamily\footnotesize,breaklines=true]
tools/call session_ctl
{"limits": {"states": 40000}, "lean": "protocol PlaceOrder ...", "formula": {"AF": "terminal"}}

answer (excerpt)
{"verdict": "proved",
 "limitsApplied":
 {"warning":
  "limits changed for this call: states 40000 (default 20000); a raised limit trades memory and time for the answer",
  "limits": {"states": {"requested": 40000, "applied": 40000}}}}
\end{lstlisting}


The operator caps what a call may ask for with \texttt{LEANFM\_LIMIT\_CEILING}, a
factor over each default (100 unless set), and a request past it is capped and says
so. Like every variable the server reads, it is checked when the server starts: a
malformed value stops the server at once rather than surfacing later in its logs.

\section{The logic tools}

The logic tools work on drafts and traces rather than protocol text: they explain a
formula, classify a requirement as observable or not, check a trace of events with
causes, and say what an actor can derive from the messages it has seen.

% Generated by LogicInteractionExample.lean from LeanFM.Logic.run; do not edit.
\medskip\noindent\textbf{A never-claim that fails (\texttt{check\_ctl})}\par\nopagebreak
\begin{lstlisting}[language={},breaklines=true]
{"task": "mcp_logic_consult",
 "formula": {"never": {"on": "LogicInputRejected"}}}
\end{lstlisting}
\noindent returns
\begin{lstlisting}[language={},breaklines=true]
{"verdict": "disproved",
 "tool": "check_ctl",
 "task": "mcp_logic_consult",
 "readingIndented":
 "on every path, never LogicInputRejected was just delivered",
 "reading":
 "on every path, never LogicInputRejected was just delivered",
 "reachableNodes": 4,
 "notation": "never LogicInputRejected",
 "mode": "never",
 "formula": {"never": {"on": "LogicInputRejected"}},
 "evidence":
 {"path":
  [{"state": "drafting", "message": null},
   {"state": "interpreting", "message": "LogicConsultRequested"},
   {"state": "rejected", "message": "LogicInputRejected"}],
  "kind": "counterexample"},
 "deadlocks": []}
\end{lstlisting}

\medskip\noindent\textbf{A requirement over hidden state (\texttt{classify\_requirement})}\par\nopagebreak
\begin{lstlisting}[language={},breaklines=true]
{"task": "mcp_logic_consult", "predicates": ["the cache is consistent"]}
\end{lstlisting}
\noindent returns
\begin{lstlisting}[language={},breaklines=true]
{"verdict": "classified",
 "tool": "classify_requirement",
 "reading":
 "some atoms are implementation-only; restate them over proposed events before the requirements gate",
 "atoms":
 [{"restated":
   {"AG":
    {"implies":
     [{"on": "CacheConsistentObserved"},
      {"field":
       {"value": "true",
        "name": "holds",
        "message": "CacheConsistentObserved"}}]}},
   "reason":
   "no declared message or state reveals this; no event trace can test it",
   "proposedEvent":
   {"to": "an observer that can check it",
    "message": "CacheConsistentObserved",
    "from": "the actor that owns this state",
    "fields": ["subject", "holds", "observed_at"]},
   "class": "implementation-only",
   "atom": "the cache is consistent",
   "advice":
   "move it to Implementation.lean as an internal invariant, or restate it over an event that exposes it"}]}
\end{lstlisting}

\medskip\noindent\textbf{Designing a sequence diagram by writing events (\texttt{check\_trace})}\par\nopagebreak
\begin{lstlisting}[language={},breaklines=true]
{"events":
 [{"task": "mcp_logic_consult",
   "src": "Agent",
   "message": "LogicConsultRequested",
   "id": "e1",
   "dst": "LeanFMServer"},
  {"task": "mcp_logic_consult",
   "src": "LeanFMServer",
   "prevIds": ["e1"],
   "message": "LogicVerdict",
   "id": "e2",
   "dst": "Agent",
   "causeIds": ["e1"]}]}
\end{lstlisting}
\noindent returns
\begin{lstlisting}[language={},breaklines=true]
{"verdict": "admissible",
 "tool": "check_trace",
 "reading": "2 events over 1 task(s); 0 concurrent pair(s)",
 "issues": [],
 "concurrentCount": 0,
 "concurrent": [],
 "acceptance":
 [{"task": "mcp_logic_consult",
   "states": ["answered"],
   "reachedTerminal": true,
   "accepted": true}]}
\end{lstlisting}
\noindent with the generated Mermaid sequence diagram
\begin{lstlisting}[language={},breaklines=true]
sequenceDiagram
    participant Agent
    participant LeanFMServer
    Agent->>LeanFMServer: e1 LogicConsultRequested
    LeanFMServer->>Agent: e2 LogicVerdict
\end{lstlisting}

\medskip\noindent\textbf{Who can derive what (\texttt{knows})}\par\nopagebreak
\begin{lstlisting}[language={},breaklines=true]
{"query": {"atom": "nb"},
 "initial": [{"sk": "I"}],
 "events":
 [{"term":
   {"enc":
    {"key": {"pk": "A"},
     "body": {"tuple": [{"atom": "na"}, {"atom": "nb"}]}}},
   "src": "Bob",
   "id": "e1",
   "dst": "Ivo"},
  {"term": {"enc": {"key": {"pk": "I"}, "body": {"atom": "nb"}}},
   "src": "Alice",
   "id": "e2",
   "dst": "Ivo",
   "causeIds": ["e1"]}],
 "at": "e2",
 "actor": "Ivo"}
\end{lstlisting}
\noindent returns
\begin{lstlisting}[language={},breaklines=true]
{"visible": ["e1", "e2"],
 "verdict": "proved",
 "tool": "knows",
 "reading": "Ivo can derive nb at e2",
 "query": "nb",
 "holds":
 [{"term": "sk(I)", "how": "known at the start", "from": "initial"},
  {"term": "{na, nb}pk(A)",
   "how": "delivered in e1 (Bob to Ivo)",
   "from": "e1"},
  {"term": "{nb}pk(I)",
   "how": "delivered in e2 (Alice to Ivo)",
   "from": "e2"},
  {"term": "nb", "how": "opened {nb}pk(I) with sk(I)", "from": "e2"}],
 "evidence":
 {"kind": "derivation",
  "derivation":
  {"term": "nb", "how": "opened {nb}pk(I) with sk(I)", "held": "e2"}},
 "causalPast": ["e1", "e2"],
 "at": "e2",
 "actor": "Ivo"}
\end{lstlisting}


